% ============================================
% DOC-22 (EN)
% English translation — SPEC-935-R355
% ============================================
\section*{\textbf{DOC-22} --- Internal Memorandum No. 47/2026}
\markboth{DOC-22 --- Internal Memorandum No. 47/2026}{}

\noindent \textit{Document extracted from the forensic file.}

\vspace{0.5em}\noindent\rule{\textwidth}{0.5pt}\vspace{0.5em}

\noindent {\footnotesize\texttt{PUBLIC ARCHIVE $\cdot$ DOCUMENTARY EVIDENCE $\cdot$ CASE 1,260}}\\[-0.2em]

\noindent \textbf{Minas Gerais Public Archive --- Documentation Directorate}\\
\noindent \textbf{To:} Special Collections Sector. \textbf{Cc:} Administrative Superintendence.\\
\noindent \textbf{Subject:} Closure of analysis of the object catalogued under record 1,260.\\
\noindent\textbf{Date:} August 1, 2026.

\noindent Dear colleagues,

\noindent 1. Whereas the notebook object attributed to record 1,260 (Barbacena Hospital-Colony) was submitted to external analysis by the forensic psychiatrist Dr. Lúcia Mendes (CRM-MG 28,391), and the said professional requested clinical leave before concluding her assessment;

\noindent 2. Whereas the re-diagnosis report (DOC-08) has already been attached to the original file and there is, at this moment, no judicial, administrative, or academic determination that justifies keeping the object in circulation;

\noindent 3. Considering the institution's document security guidelines and the principle of collection stability;

\noindent \textbf{I hereby determine:}

\noindent a) The immediate closure of public consultation of object 1,260;

\noindent b) The storage of the object in a definitive vault, with no provision for new access;

\noindent c) The notification of all involved parties that the case is closed and no new handling requests will be admitted.

\noindent \textbf{Rationale:} the object does not constitute material evidence of a crime, is not part of any ongoing legal action, and does not qualify as a document of exceptional historical value that would justify risk to the technical staff. The continuation of its circulation endangers document security and the psychological integrity of the professionals involved.

\noindent \textbf{Additional recommendation:} I suggest that the directorate assess whether it is advisable not to register this memorandum in the computerized document management system. Computerized registration makes the document locatable. And a locatable document can be requested again.

\noindent \textit{Ronaldo Torres}\\
\noindent \textit{Minas Gerais Public Archive --- Director of Documentation}
\noindent \textit{Handwritten annotation in the top margin, blue ink, cursive handwriting:}
\noindent \textquotedblleft{}\textit{Ronaldo --- it was not possible. The object was removed from the vault the morning after the memorandum. The tracking system records the vault opening at 3:14 a.m. No one was identified on the cameras. The door was sealed. The seal was intact. Matter filed as requested.\textquotedblright{}}

\noindent \textquotedblleft{}\textit{--- Special Collections Sector, 08/02/2026.\textquotedblright{}}

\vspace{0.5em}
\noindent\textit{\textbf{↪ Routes:} \rota{DOC-01} \rotasep \rota{LUC-02} \rotasep \rota{DOC-23}}
